\documentclass[12pt,a4paper]{article}

% ---------------------------------------------------------
% Idioma y tipografía
% ---------------------------------------------------------
\usepackage[utf8]{inputenc}
\usepackage[T1]{fontenc}
\usepackage[provide=*,spanish]{babel}
\usepackage{mathpazo}
\usepackage{microtype}

% ---------------------------------------------------------
% Matemática y formato
% ---------------------------------------------------------
\usepackage{amsmath,amssymb,amsthm}
\usepackage{geometry}
\usepackage{enumitem}
\usepackage[hidelinks]{hyperref}

\geometry{
	margin=2.5cm
}

\newcommand{\R}{\mathbb{R}}
\newcommand{\dd}{\mathrm{d}}

\title{Ecuaciones Diferenciales Ordinarias\\Notas para el docente: diferenciales y formas diferenciales}
\author{Benjamín Marcolongo}
\date{}

\begin{document}

	\maketitle

	\section{Propósito}

	Estas notas desarrollan una interpretación geométrica de la forma diferencial de una EDO de primer orden. El objetivo es precisar qué representan los símbolos \(dx\) y \(dy\), por qué una expresión
	\[
	M(x,y)\,dx+N(x,y)\,dy=0
	\]
define direcciones tangentes y cómo se relacionan las ecuaciones de Pfaff con las soluciones implícitas, las ecuaciones exactas, los factores integrantes y la posible pérdida de soluciones durante una manipulación algebraica.

	La notación elemental sugiere que \(dx\) y \(dy\) son incrementos infinitesimales. Esa intuición es útil, pero la formulación geométrica es más precisa: son secciones del fibrado cotangente y actúan linealmente sobre vectores tangentes.

	\clearpage
	\section{El diferencial como covector}

	Sea \(M\) una variedad suave y sea
	\[
	f\colon M\longrightarrow\R
	\]
una función suave. El diferencial de \(f\) en \(p\in M\) es el covector
	\[
	df_p\in T_p^*M
	\]
definido por
	\[
	df_p(v)=v[f],
	\]
para todo \(v\in T_pM\). Si \(\gamma(s)\) es una curva con \(\gamma(0)=p\) y \(\dot{\gamma}(0)=v\), entonces
	\[
	df_p(v)
	=
	\left.\frac{d}{ds}f\bigl(\gamma(s)\bigr)\right|_{s=0}.
	\]

	En coordenadas locales \((x^1,\dots,x^n)\),
	\[
	df
	=
	\sum_{i=1}^{n}
	\frac{\partial f}{\partial x^i}\,dx^i.
	\]
Los objetos \(dx^i\) son los diferenciales de las funciones coordenadas
	\[
	x^i\colon M\longrightarrow\R.
	\]
Forman la co-base dual de \(\{\partial/\partial x^i\}\):
	\[
	dx^i\left(\frac{\partial}{\partial x^j}\right)
	=
	\delta^i_j.
	\]

	En \(\R^2\), si
	\[
	v=v^x\frac{\partial}{\partial x}
	+v^y\frac{\partial}{\partial y},
	\]
entonces
	\[
	dx(v)=v^x,
	\qquad
	dy(v)=v^y.
	\]
Así, \(dx\) y \(dy\) no son vectores ni números pequeños: son funcionales lineales que extraen las componentes de un vector tangente.

	\subsection{Relación con la linealización}

	En un espacio vectorial, el diferencial es la parte lineal principal del incremento:
	\[
	f(p+v)
	=
	f(p)+df_p(v)+o(\lVert v\rVert).
	\]
Esta fórmula justifica la intuición elemental. El incremento finito \(\Delta f\) solo coincide aproximadamente con \(df\), mientras que \(df_p(v)\) está definido exactamente como la linealización de \(f\) en \(p\).

	\clearpage
	\section{Pullback a una curva}

	Sea
	\[
	\gamma\colon I\longrightarrow M,
	\qquad
	s\longmapsto \gamma(s),
	\]
una curva suave. Si \(\alpha\) es una 1-forma sobre \(M\), su pullback es la 1-forma sobre \(I\)
	\[
	\gamma^*\alpha
	=
	\alpha_{\gamma(s)}\bigl(\dot{\gamma}(s)\bigr)\,ds.
	\]

	Para una curva del plano
	\[
	\gamma(s)=\bigl(x(s),y(s)\bigr),
	\]
se obtiene
	\[
	\gamma^*(dx)=\dot{x}(s)\,ds,
	\qquad
	\gamma^*(dy)=\dot{y}(s)\,ds.
	\]
Esta es la interpretación precisa de las expresiones simbólicas
	\[
	dx=\dot{x}\,ds,
	\qquad
	dy=\dot{y}\,ds
	\]
cuando se restringen los diferenciales a una curva parametrizada.

	Si se elige \(x\) como parámetro, de modo que \(\gamma(x)=(x,y(x))\), entonces
	\[
	\gamma^*(dx)=dx,
	\qquad
	\gamma^*(dy)=\dot{y}\,dx,
	\]
donde el punto representa en este caso derivación respecto de \(x\).

	Esta elección requiere que \(x\) sea una coordenada válida a lo largo de la curva, es decir, que \(\dot{x}\neq 0\) en la parametrización original. Las tangentes verticales son precisamente los puntos en los que esta descripción deja de ser posible.

	\clearpage
	\section{La ecuación de Pfaff}

	Sea
	\[
	\alpha
	=
	M(x,y)\,dx+N(x,y)\,dy
	\]
una 1-forma sobre una región \(U\subseteq\R^2\). La ecuación de Pfaff
	\[
	\alpha=0
	\]
no afirma que la 1-forma sea idénticamente nula. Indica que una curva solución \(\gamma\) debe satisfacer
	\[
	\gamma^*\alpha=0,
	\]
o equivalentemente,
	\[
	\alpha_{\gamma(s)}\bigl(\dot{\gamma}(s)\bigr)=0.
	\]
En coordenadas,
	\[
	M\bigl(x(s),y(s)\bigr)\dot{x}(s)
	+
	N\bigl(x(s),y(s)\bigr)\dot{y}(s)
	=0.
	\]

	En cada punto donde \(\alpha_p\neq 0\), la forma define la distribución unidimensional
	\[
	\mathcal{D}_p=\ker\alpha_p\subset T_pU.
	\]
Las curvas solución son las curvas integrales de esta distribución. Un campo vectorial que genera localmente el núcleo es
	\[
	X
	=
	N\frac{\partial}{\partial x}
	-
	M\frac{\partial}{\partial y},
	\]
pues
	\[
	\alpha(X)=MN-NM=0.
	\]

	Si se usa \(x\) como parámetro y \(N\neq 0\), la ecuación se reduce a
	\[
	M(x,y)+N(x,y)\dot{y}=0,
	\]
es decir,
	\[
	\dot{y}=-\frac{M(x,y)}{N(x,y)}.
	\]
La forma normal no es entonces otra ecuación: es la representación de la misma distribución en una carta y con una elección particular de parámetro.

	\subsection{Reparametrización}

	La condición \(\gamma^*\alpha=0\) es invariante ante reparametrizaciones regulares. Si \(\widetilde{\gamma}=\gamma\circ\sigma\), el vector tangente se multiplica por \(\dot{\sigma}\), pero la anulación de \(\alpha(\dot{\gamma})\) no cambia. Esta invariancia queda parcialmente oculta al escribir directamente \(\dot{y}=f(x,y)\).

	\clearpage
	\section{Interpretación geométrica en el plano}

	Usando el producto euclídeo para identificar covectores y vectores, la 1-forma
	\[
	\alpha=M\,dx+N\,dy
	\]
corresponde al vector normal
	\[
	\alpha^\sharp=(M,N).
	\]
La condición
	\[
	\alpha(\dot{\gamma})=0
	\]
expresa que el vector tangente es perpendicular a \((M,N)\).

	El ejemplo básico es
	\[
	\alpha=x\,dx+y\,dy.
	\]
Como
	\[
	\alpha=d\left(\frac{x^2+y^2}{2}\right),
	\]
las curvas integrales satisfacen
	\[
	x^2+y^2=C.
	\]
Geométricamente, \((x,y)\) es normal a las circunferencias, mientras que
	\[
	X
	=
	y\frac{\partial}{\partial x}
	-
	x\frac{\partial}{\partial y}
	\]
es tangente a ellas.

	Esta identificación euclídea es útil en el plano, pero no es intrínseca: transformar una 1-forma en un vector requiere una métrica. La formulación \(\ker\alpha\) no necesita elegir ninguna.

	\clearpage
	\section{Formas exactas y primeras integrales}

	La ecuación es exacta si existe una función \(F\colon U\to\R\) tal que
	\[
	\alpha=dF.
	\]
Para una curva integral,
	\[
	0
	=
	\gamma^*\alpha
	=
	\gamma^*(dF)
	=
	d(F\circ\gamma).
	\]
Por lo tanto,
	\[
	F\circ\gamma=C,
	\]
y las soluciones son curvas de nivel
	\[
	F(x,y)=C.
	\]

	Si
	\[
	\alpha=M\,dx+N\,dy,
	\]
entonces
	\[
	d\alpha
	=
	\left(
	\frac{\partial N}{\partial x}
	-
	\frac{\partial M}{\partial y}
	\right)
	dx\wedge dy.
	\]
Toda forma exacta es cerrada. En una región simplemente conexa del plano, una 1-forma cerrada es exacta. Localmente, esta equivalencia es una consecuencia del lema de Poincaré.

	\subsection{Soluciones implícitas y valores regulares}

	Si \(dF_p\neq 0\), el conjunto de nivel \(F^{-1}(C)\) es localmente una subvariedad unidimensional. Para representarlo como una gráfica \(y=y(x)\) se necesita además
	\[
	\frac{\partial F}{\partial y}(p)\neq 0.
	\]
En ese caso,
	\[
	\dot{y}
	=
	-\frac{\partial F/\partial x}{\partial F/\partial y}.
	\]
Si \(dF_p=0\), el conjunto de nivel puede dejar de ser regular, bifurcarse o reducirse a un punto. Esta singularidad del conjunto de nivel no debe identificarse automáticamente con una solución singular de la EDO.

	\subsection{Una obstrucción global}

	En \(\R^2\setminus\{0\}\), la forma
	\[
	\alpha
	=
	\frac{-y\,dx+x\,dy}{x^2+y^2}
	\]
es cerrada, pero no exacta globalmente. Localmente es \(d\theta\), mientras que sobre una circunferencia orientada positivamente
	\[
	\oint \alpha=2\pi.
	\]
El ejemplo muestra que la condición \(d\alpha=0\) garantiza potenciales locales, pero la topología del dominio puede impedir la existencia de una primera integral global univaluada.

	\clearpage
	\section{Integrabilidad y factores integrantes}

	En una variedad de dimensión arbitraria, la distribución de hiperplanos \(\ker\alpha\) es integrable si satisface la condición de Frobenius
	\[
	\alpha\wedge d\alpha=0.
	\]
En dimensión dos esta condición se cumple automáticamente, porque \(\alpha\wedge d\alpha\) sería una 3-forma. Por lo tanto, toda 1-forma suave y no nula sobre una superficie define localmente una foliación por curvas.

	Esto no significa que \(\alpha\) sea exacta. Significa que localmente existe una función \(F\) y una función no nula \(g\) tales que
	\[
	\alpha=g\,dF.
	\]
Equivalentemente, existe un factor integrante \(\mu=1/g\) para el cual
	\[
	\mu\alpha=dF.
	\]

	Como ejemplo, consideremos
	\[
	\alpha=dy-y\,dx.
	\]
No es cerrada:
	\[
	d\alpha=dx\wedge dy.
	\]
Sin embargo, el factor
	\[
	\mu(x)=e^{-x}
	\]
produce
	\[
	e^{-x}\alpha
	=
	e^{-x}\,dy-ye^{-x}\,dx
	=
	d\left(ye^{-x}\right).
	\]
Por consiguiente,
	\[
	ye^{-x}=C,
	\]
o
	\[
	y=Ce^x.
	\]

	Multiplicar \(\alpha\) por una función \(\mu\) que no se anula preserva la distribución:
	\[
	\ker(\mu\alpha)=\ker\alpha.
	\]
Si \(\mu\) se anula, esta equivalencia falla porque el rango de la forma puede cambiar sobre el conjunto de ceros.

	\clearpage
	\section{Separación de variables y dominio}

	La ecuación
	\[
	\dot{y}=g(x)h(y)
	\]
puede asociarse con la forma
	\[
	\alpha=dy-g(x)h(y)\,dx.
	\]
Sobre una región donde \(h(y)\neq 0\), multiplicar por \(1/h(y)\) da
	\[
	\widetilde{\alpha}
	=
	\frac{1}{h(y)}\,dy-g(x)\,dx.
	\]
Esta forma es exacta:
	\[
	\widetilde{\alpha}
	=
	d\left(
	\int^y\frac{d\eta}{h(\eta)}
	-
	\int^x g(\xi)\,d\xi
	\right).
	\]
La separación de variables puede interpretarse así como la búsqueda inmediata de un factor integrante que convierte la ecuación de Pfaff en una forma exacta.

	La restricción \(h(y)\neq 0\) no es un detalle algebraico. Si \(h(y_\ast)=0\), la curva
	\[
	y=y_\ast
	\]
puede ser una solución integral de la forma original, pero queda fuera del dominio de \(\widetilde{\alpha}\).

	Para
	\[
	\dot{y}=x\sqrt{y},
	\qquad y\geq 0,
	\]
la forma original
	\[
	\alpha=dy-x\sqrt{y}\,dx
	\]
se anula sobre el vector tangente a la curva \(y=0\). En cambio,
	\[
	\frac{dy}{\sqrt{y}}-x\,dx
	\]
solo está definida para \(y>0\). La solución nula no desaparece por un misterio analítico: fue eliminada al cambiar el dominio de la 1-forma.

	Además, \(\sqrt{y}\) no es localmente Lipschitz en \(y=0\). La degeneración que permite pegar ramas no nulas con segmentos de la solución nula es consistente con la pérdida de unicidad del campo de direcciones.

	\clearpage
	\section{Puntos singulares de la forma de Pfaff}

	La discusión anterior presupone \(\alpha_p\neq 0\). Si
	\[
	M(p)=N(p)=0,
	\]
entonces
	\[
	\ker\alpha_p=T_pU,
	\]
y la distribución deja de tener rango constante. En ese punto, la ecuación no selecciona una única dirección tangente.

	Conviene distinguir tres nociones:

	\begin{enumerate}[label=\arabic*.]
		\item un punto singular de la 1-forma, donde \(\alpha_p=0\);
		\item un punto crítico de una primera integral, donde \(dF_p=0\);
		\item una solución singular, entendida como una solución que no pertenece a una familia uniparamétrica dada.
	\end{enumerate}

	Estas nociones pueden relacionarse en ejemplos particulares, pero no son equivalentes.

	La formulación geométrica natural de una ecuación de Clairaut no vive solamente en el plano \((x,y)\), sino en el espacio de jets \(J^1(\mathbb{R},\mathbb{R})\), con coordenadas \((x,y,p)\) y forma de contacto
	\[
	\vartheta=dy-p\,dx.
	\]
Para la ecuación
	\[
	y=xp+p^2,
	\]
las soluciones corresponden a curvas integrales de \(\ker\vartheta\) contenidas en esa superficie. La familia regular se obtiene fijando \(p=C\). La envolvente aparece donde la proyección al plano \((x,y)\) se vuelve singular respecto del parámetro:
	\[
	\frac{\partial}{\partial p}\left(y-xp-p^2\right)=-(x+2p)=0.
	\]
De aquí,
	\[
	p=-\frac{x}{2},
	\qquad
	y=-\frac{x^2}{4}.
	\]
Así, la singularidad relevante es una singularidad de la proyección de la familia legendriana, no simplemente un cero de una 1-forma sobre el plano.

	\clearpage
	\section{Puente pedagógico}

	Para una primera clase no introduciría fibrados cotangentes ni pullbacks. Presentaría dos niveles de lectura.

	\subsection{Lectura elemental}

	Los símbolos \(dx\) y \(dy\) representan las componentes horizontal y vertical de un desplazamiento tangente. La ecuación
	\[
	M\,dx+N\,dy=0
	\]
selecciona los desplazamientos permitidos. En el plano,
	\[
	(M,N)\cdot(dx,dy)=0,
	\]
por lo que la dirección tangente es perpendicular a \((M,N)\).

	\subsection{Lectura matemática subyacente}

	En realidad, \(dx\) y \(dy\) son la co-base coordenada, y la ecuación selecciona el núcleo de una 1-forma. La escritura con desplazamientos infinitesimales funciona porque el pullback a una curva convierte
	\[
	dx\mapsto\dot{x}\,ds,
	\qquad
	dy\mapsto\dot{y}\,ds.
	\]

	La frase que utilizaría con los alumnos es:

	\begin{quote}
		La forma diferencial no determina directamente la altura de la solución. Determina qué direcciones tangentes están permitidas en cada punto; las curvas que siguen esas direcciones son las soluciones.
	\end{quote}

	El ejemplo
	\[
	x\,dx+y\,dy=0
	\]
es particularmente eficiente porque conecta en una sola figura la dirección tangente, el vector normal, el diferencial exacto y la solución implícita
	\[
	x^2+y^2=C.
	\]

	\clearpage
	\section{Resumen conceptual}

	\begin{itemize}[leftmargin=*]
		\item El símbolo \(dx\) es el diferencial de la función coordenada \(x\), es decir, un covector de la co-base coordenada.
		\item La igualdad simbólica \(dx=\dot{x}\,ds\) debe entenderse como una igualdad después de aplicar el pullback a una curva.
		\item Una ecuación de Pfaff \(\alpha=0\) define la distribución \(\ker\alpha\); sus soluciones son las curvas integrales de esa distribución.
		\item La forma normal \(\dot{y}=f(x,y)\) requiere elegir \(x\) como parámetro y excluye localmente tangentes verticales.
		\item Si \(\alpha=dF\), las soluciones son curvas de nivel de una primera integral \(F\).
		\item En dimensión dos, toda 1-forma no nula es localmente integrable, aunque no necesariamente exacta.
		\item Dividir una forma por una función que puede anularse cambia su dominio y puede eliminar soluciones.
		\item Los puntos singulares de una 1-forma, los puntos críticos de una primera integral y las soluciones singulares son conceptos diferentes.
	\end{itemize}

	\section*{Bibliografía}

	\begin{itemize}[leftmargin=*]
		\item Lee, John M. \textit{Introduction to Smooth Manifolds}.
		\item Spivak, Michael. \textit{Calculus on Manifolds}.
		\item Arnold, Vladimir I. \textit{Ordinary Differential Equations}.
		\item Zill, Dennis G. \textit{A First Course in Differential Equations with Modeling Applications}.
	\end{itemize}

	\subsection*{Nota sobre la elaboración del material}

	Este material fue elaborado por Benjamín Marcolongo con la asistencia de \textbf{GPT-5.6 Sol}, un modelo de OpenAI utilizado a través del entorno Codex, para la organización, redacción y revisión del contenido.

\end{document}
